\documentclass[11pt]{article}

\usepackage[margin=1in]{geometry}
\usepackage{amsmath,amssymb,amsthm,mathtools}
\usepackage{enumitem}
\usepackage{booktabs}
\usepackage{tabularx}
\usepackage{microtype}
\usepackage[hidelinks]{hyperref}
\usepackage{xurl}

\newtheorem{definition}{Definition}
\newtheorem{lemma}{Lemma}
\newtheorem{remark}{Remark}

\title{Challenge Branch Maps for Successor Maintenance:\\[0.5ex]
\large Piece-Keyed Escalation from Audience Routes to Approved Terminal Fields}
\author{}
\date{Unified archive note v0.03 (2026-03-21; archive v488)}

\begin{document}
\maketitle

\begin{abstract}
Audience-keyed challenge routes answer where a given reader should begin.
Question-keyed stop profiles answer which semantic pieces exist and which narrow fields are allowed to discharge each piece.
One operational layer is still missing if those two objects are left alone: once a specific semantic piece is actually contested, what exact escalation walk is allowed from that audience route to one approved terminal field?
This note isolates that bridge.
It defines a piece-keyed \emph{challenge branch map} that records the exact route-plus-piece walk, keeping coarse audience ladders separate from exact piece-local escalation.
The result is a compact, auditable home for contested-piece traces without overloading either routes or stop profiles.
\end{abstract}

\paragraph{Non-redundancy note.}
Synthesis~31 owns the downstream normal form, Synthesis~36 owns audience-keyed start objects and coarse escalation order, Synthesis~37 owns question families and approved terminal fields, Synthesis~54 owns the answer-side terminal citation normal form, Synthesis~74 owns the paired terminal citation discipline, Synthesis~75 owns the paired cutover rule, and Synthesis~17 owns the worked instantiation.
This note owns only the piece-keyed escalation bridge from one route to one approved terminal field. Later notes should reopen it only when the exact piece-keyed escalation walk is itself live.

\paragraph{Scope.}
This is not a new receipt object, verdict, or maintenance stage.
It is a publication-interface note about what exact escalation walk is allowed once one semantic piece of a downstream question is actually challenged.

\paragraph{Imports.}
Audience-keyed routes are imported from Synthesis~36.\cite{series:challengeroutes}
Question-keyed stop semantics are imported from Synthesis~37.\cite{series:challengestops}
The concrete worked instantiation is imported from Synthesis~17.\cite{series:workedexample}

\paragraph{Exports.}
This note exports (i) piece-keyed challenge branches, (ii) a route-link rule saying every branch begins at a declared audience-route start object, (iii) a piece-link rule saying every branch terminates at an approved terminal field from the companion stop profile, and (iv) a local-reopen rule saying that later terse papers should reopen this note only when the exact piece-keyed escalation walk is itself live. Otherwise later terse papers should default to the answer-side terminal citation normal form in Synthesis~54 together with the paired terminal discipline and paired cutover rule in Synthesis~74--75. Exact question-family stop semantics remain outside scope here and belong to Synthesis~37.\cite{series:challengestops,series:challengeanswerreviewmenus,series:pairedterminalcitationdiscipline,series:pairedterminalcutoverrules}

\section{Why the branch layer should be separate}
A release-note route can say ``start at the outward-facing sentence and escalate if challenged.''
A stop profile can say ``the continuity half of that sentence may stop at the continuity verdict.''
Neither statement alone yet fixes the exact contested-piece walk.
Without a branch layer, later notes tend to improvise that walk ad hoc from a coarse audience ladder, and the resulting proof traces drift.
So the piece-specific escalation bridge deserves one small home of its own.

\begin{definition}[Challenge branch]
Fix one concrete base$\to$successor comparison, one audience-keyed route $r$, and one contested semantic piece $p$ approved by the companion stop profile for $r$'s question family.
A \emph{challenge branch}
\[
B(r,p)=(s_0,s_1,\dots,s_m; t)
\]
consists of a finite nonempty sequence of already-owned artifacts or artifact slices beginning at the route's declared start object $s_0$ and ending at a terminal field $t$ that is approved for piece $p$ by the companion stop profile.
\end{definition}

\begin{definition}[Route-link rule]
A challenge branch satisfies the \emph{route-link rule} if its first step agrees with the start object declared by its companion audience-keyed route.
\end{definition}

\begin{definition}[Piece-link rule]
A challenge branch satisfies the \emph{piece-link rule} if its terminal field is one of the approved terminal fields for the same piece key in the companion stop profile.
\end{definition}

\begin{lemma}[Piece-keyed branching is audit-safe]
Assume every challenge branch for a concrete successor comparison satisfies the route-link rule and the piece-link rule.
Then later papers may record exact contested-piece escalation walks without changing audience navigation, stop semantics, or ownership of narrow terminal fields.
\end{lemma}

\begin{proof}
The route-link rule prevents a branch map from inventing a new audience entry point.
The piece-link rule prevents it from inventing a new terminal field.
So the branch layer can refine one coarse audience ladder into one exact piece-local walk while reusing the same already-owned start objects and terminal fields.
It adds operational precision, not new semantics.
\end{proof}

\begin{table}[t]
\centering
\small
\begin{tabularx}{\linewidth}{@{}l X X@{}}
\toprule
Question / audience slice & Contested piece & Typical branch end \\
\midrule
Release-note public status & continuity piece & continuity verdict token \\
Release-note public status & closure piece & closure-status token \\
Referee package completeness & role-row piece & closure-ledger rows \\
Auditor public-status token & token piece & status-envelope token \\
Auditor compact diff & recompute piece & verifier recomputation slice \\
\bottomrule
\end{tabularx}
\caption{Challenge branches are route-plus-piece traces. They are finer than audience ladders and narrower than wrapper prose.}
\label{tab:branchmap}
\end{table}

\section{Worked example pointer}
The maintained worked bundle now ships \path{example_successor_challenge_branches.json}.\cite{series:workedexample}
It records fifteen branch entries for the canned Case-B successor.
Each entry binds one audience route and one contested semantic piece to an exact escalation sequence and one approved terminal field.
For example, the release-note public-status route splits into distinct continuity and closure branches, while the referee package-completeness route splits into a role-row branch and a closure-token branch.
That keeps the route note coarse, the stop-profile note semantic, and the exact contested-piece trace in one explicit object.\cite{series:challengeroutes,series:challengestops}

\paragraph{Why this helps the archive.}
The series can now keep audience ladders, stop semantics, and exact piece-local escalation separate.
That makes the worked example more auditable without lengthening the public wrapper sentence or forcing every downstream note to improvise a contested-piece proof trace.


\paragraph{A minimal public challenge-branch card.}
\begin{table}[t]
\centering
\small
\begin{tabularx}{\linewidth}{@{}lX@{}}
\toprule
Field & Minimal public answer \\
\midrule
Release-note continuity branch & \nolinkurl{release_note_public_status_sentence_continuity_piece} starts at \nolinkurl{example_successor_lineage_notice.json}, steps through \nolinkurl{example_successor_claim_support_map.json}, ends at \nolinkurl{example_successor_continuity_verdict.json}, and stops at the approved field \nolinkurl{continuity_decision} \\
Referee package-completeness branch & \nolinkurl{referee_package_completeness_role_rows_piece} starts and ends at \nolinkurl{example_release_closure_ledger.json} and stops at the approved field \nolinkurl{role_rows} \\
Release-note recompute branch & \nolinkurl{release_note_compact_diff_recompute_piece} starts at \nolinkurl{example_compare_report.json}, steps through \nolinkurl{example_verifier_report.json}, and stops at the approved field \nolinkurl{recomputed} \\
Selection rule & Stop here only while the live issue is the exact contested-piece walk from one maintained audience route to one approved terminal field; do not widen to family-level coverage, exact surface realization, or exact terminal ownership unless those neighboring owners have become the honest moving part \\
Left/right boundary & Reopen Synthesis~36 when the live issue is the audience-keyed start object or coarse escalation order, reopen Synthesis~37 when the contested semantic piece or approved stop field changes, and widen to Synthesis~39--41 once the live issue becomes coverage, exact surface binding, or exact narrow owner rather than the branch walk itself \\
Paired-terminal boundary & Once one maintained branch still forces the same exact contested-piece walk, later terse papers should usually stop at Synthesis~54 together with Synthesis~74--75 rather than reopen this note merely to restate a stable answer-side handoff or paired cutover judgment \\
\bottomrule
\end{tabularx}
\caption{A minimal public challenge-branch card. This is the smallest exact escalation object later terse papers can quote directly when the live issue is how one contested semantic piece travels from one maintained audience route to one approved terminal field.}
\label{tab:minimal-public-challenge-branch-card}
\end{table}

This card is intentionally non-decorative: it pins the actual worked branch keys, path sequences, and terminal stop fields without silently re-owning the earlier audience-route or stop-profile owners or the later coverage / surface-hook / terminal-owner layers that neighboring notes already own.\cite{series:workedexample}

\paragraph{A forced public challenge-branch matrix.}
\begin{table}[t]
\centering
\small
\begin{tabularx}{\linewidth}{@{}lX@{}}
\toprule
Field & Forced public answer \\
\midrule
Release-note continuity floor & \nolinkurl{release_note_public_status_sentence_continuity_piece} still forces the same branch only while the same route key and piece key agree, the same start path remains \nolinkurl{example_successor_lineage_notice.json}, the same intermediate path \nolinkurl{example_successor_claim_support_map.json} remains in place, and the same approved stop field \nolinkurl{continuity_decision} remains at \nolinkurl{example_successor_continuity_verdict.json} \\
Referee package-completeness floor & \nolinkurl{referee_package_completeness_role_rows_piece} still forces the same branch only while the same route key and piece key agree, the same start path remains \nolinkurl{example_release_closure_ledger.json}, the walk remains one-step, and the same approved stop field \nolinkurl{role_rows} remains at that ledger \\
Release-note recompute floor & \nolinkurl{release_note_compact_diff_recompute_piece} still forces the same branch only while the same route key and piece key agree, the same start path remains \nolinkurl{example_compare_report.json}, the same intermediate path remains \nolinkurl{example_verifier_report.json}, and the same approved stop field \nolinkurl{recomputed} remains at that verifier report \\
Route/stop boundary floor & Those shortcuts remain honest only while the companion audience route still names the same declared start object and the companion stop profile still approves the same terminal field for the same contested piece; the branch does not get to keep the same answer once either neighboring owner changes under it \\
Staleness trigger floor & This matrix goes stale as soon as the route key changes, the piece key changes, a path is added / removed / reordered, the declared start object changes, the terminal stop field changes, the terminal path changes, or the maintained branch ceases to be the same exact contested-piece walk \\
Escalation pointer & Stop here only while one maintained route-plus-piece pair still forces the same exact contested-piece walk; reopen Synthesis~36 when the audience route moves, reopen Synthesis~37 when the semantic piece or approved stop field moves, and widen to Synthesis~39--41 when the honest issue is coverage, exact surface realization, or exact terminal ownership rather than the branch path itself \\
\bottomrule
\end{tabularx}
\caption{A forced public challenge-branch matrix. This is the smallest exact escalation object later terse papers can quote directly when the live issue is whether one maintained route-plus-piece pair still forces the same contested-piece walk rather than merely which branch exists in the abstract.}
\label{tab:forced-public-challenge-branch-matrix}
\end{table}

This matrix is intentionally non-decorative: it pins the actual worked branch keys, path sequences, and approved stop fields that still force one maintained exact contested-piece walk, together with the conditions under which that forcing remains honest or goes stale, without silently re-owning the earlier route / stop-profile owners or the later coverage / surface-hook / terminal-owner layers.\cite{series:workedexample}

\paragraph{One tiny challenge-branch drill.}
For the worked release-note compact-diff question, \nolinkurl{release_note_compact_diff_recompute_piece} still forces the same exact walk from \nolinkurl{example_compare_report.json} to \nolinkurl{example_verifier_report.json} and stops at \nolinkurl{recomputed}. If the live dispute becomes ``which semantic piece is even being asked?'', the honest stop is no longer the branch map and a later terse paper must reopen Synthesis~37. If the live dispute instead stays piece-local but the verifier witness is replaced or the path order changes, then the branch matrix has gone stale and this note must be reopened rather than quoted as if the old exact contested-piece walk still held.\cite{series:workedexample}

\begin{thebibliography}{99}\small

\bibitem{series:challengeroutes}
Anonymous.
\newblock Challenge Routes and Audience Profiles for Successor Maintenance: Audience-Keyed Start Objects and Escalation Order.
\newblock Series synthesis note (Synthesis~36), 2026. (This archive.)

\bibitem{series:challengestops}
Anonymous.
\newblock Challenge Stop Profiles for Successor Maintenance: Question Families, Semantic Pieces, and Approved Terminal Fields.
\newblock Series synthesis note (Synthesis~37), 2026. (This archive.)

\bibitem{series:challengeanswerreviewmenus}
Anonymous.
\newblock Challenge Answer Review Menus for Successor Maintenance: Answer-Side Terminal Citation Normal Form.
\newblock Series synthesis note (Synthesis~54), 2026. (This archive.)

\bibitem{series:pairedterminalcitationdiscipline}
Anonymous.
\newblock Paired Terminal Citation Discipline for Review Spines: Stable Answer Stops and Adjacent Request Widening.
\newblock Series synthesis note (Synthesis~74), 2026. (This archive.)

\bibitem{series:pairedterminalcutoverrules}
Anonymous.
\newblock Paired Terminal Cutover Rules for Review Spines: When Later Terse Papers Stop, Widen, or Reopen.
\newblock Series synthesis note (Synthesis~75), 2026. (This archive.)

\bibitem{series:workedexample}
Anonymous.
\newblock Worked Example for Anonymous-DHT Receipts: Minimal Public Surface, Cross-Release Diff, and Successor Interlock.
\newblock Series synthesis note (Synthesis~17), 2026. (This archive.)

\end{thebibliography}
\end{document}
